\section{In Summary}
\label{sec:summary}

This paper asked how a research team should spend its scarce
writing hours.
Newell supplied the answer.
Know what you are good at, then aim those strengths at your
field's open problems.
Section 1 turned that advice into a loop of fourteen small
steps.
Sections 2 to 4 ran the loop's reading passes on this paper's
own authors.
Those sections computed a repertory grid from our corpus
(Figure~\ref{fig:focus}), mapped the grid to the nine ICSE'27
areas, and reviewed the literature of the territory that map
nominated.
Section 5 gave the writing rules, namely a grammar for each part
of a paper, a worked sample beside each grammar, and the rewrite
discipline that turns raw LLM output into prose.
Section 6 closed the loop with critics: a reviewer-2 prompt
built from the venue's own five review criteria, and a triage
prompt that splits the review into machine fixes and human
decisions.

We stress that none of this is the definitive way to write a
paper.
It is one way, and its parts are deliberately easy to change.
Every structure proposed here is a small grammar in a figure.
To disagree with one, edit its productions and ask an LLM to
regenerate the sample; the cost of the experiment is minutes.
Your venue may open with related work, ban the roadmap
paragraph, or demand a structured abstract.
Adjust the grammar, not your patience.

Hence our call to action.
Future work should assemble a catalog of paper types, each with
its own grammar: the empirical study, the systematic review, the
tool paper, the experience report, the vision paper.
Communities curate datasets and baselines; they could curate
paper shapes the same way, so that every writer starts from a
tested structure rather than a blank page.
Schr\"odinger asked scientists to tell everyone what they have
been doing.
A shared catalog of shapes would make that telling cheaper, and
perhaps better.

\appendix
